% ============================================================================
% Fichier  : partie3/P03_C10_cadre_juridique_deontologie.tex
% Titre    : Le Cadre Camerounais et Africain, la Déontologie
% Auteur   : MINKA MI NGUIDJOÏ Thierry Emmanuel
% Version  : 1.0 -- novembre 2026
% Corrigés : corriges/Corr_C10.tex
% ============================================================================
% Sources et degré de vérification (repris d'ACSI, chapitre 12, et Ann. D) :
%  - Loi n° 2010/012 : texte français officiel lu ; cités textuellement ici :
%    art. 7 al. 2 et 14 ; art. 61 al. 1 ; art. 4 (9), 13 al. 1 et 32 renvoyés
%    au chapitre 6 ; art. 26, 29 et 30 résumés ; art. 61 al. 3 paraphrasé.
%  - Règlement COBAC R-2016/04 : art. 1 et 2 relevés sur une base en ligne
%    (cours GCI), À COLLATIONNER avec la version du SG-COBAC.
%  - Convention de Malabo (UA, 2014) : entrée en vigueur en 2023 ; statut
%    camerounais À VÉRIFIER.
%  - Loi n° 2024/017 : présentée sans citation, À VÉRIFIER au Journal
%    officiel ; étudiée dans l'UE MSN-2S3-JUR01.
%  - Déontologie : principes éthiques des normes mondiales d'audit interne
%    de l'IIA (2024) ; menaces sur l'indépendance d'après le cadre
%    conceptuel du code de l'IESBA.
% Faits du cas : ANPP (M. ESSOMBA), BANQUE SANAGA (COBAC, rapport après
% l'incident), LiZbiZ (AUDIX CONSEIL, M. BEYALA, PHARMACAM).
% Faits nouveaux (reportés dans la bible) : en 2028, l'ANTIC annonce un
% audit de sécurité à l'ANPP ; cinq obligations légales cotées 0/3/2 ; la
% COBAC demande en 2026 un rapport sur l'incident de septembre ; M. ETEME
% demande à AUDIX d'adoucir la synthèse destinée à PHARMACAM ; la DSI de la
% SANAGA propose un « petit-déjeuner » au chef de mission de CERTIQUAL.
% ============================================================================

\chapter{Le Cadre Camerounais et Africain, la Déontologie}
\label{chap:droit}
\chaptermark{Cadre juridique et déontologie}

\epigraph{%
    « L'auditeur constate un écart au critère.\\
    Le juge, seul, qualifie l'infraction. »%
}{Règle de prudence}

\niveauChapitre{Fondamental}{%
    Maîtrise du Chapitre~\ref{chap:audit}%
}

\begin{boxobjectifs}
\begin{itemize}
    \item Distinguer, dans un texte juridique, ce qui est vérifié de ce qui
          reste à vérifier, et ne citer que le premier.
    \item Identifier les obligations d'audit et de sécurité de la loi
          n°~2010/012, et le rôle de l'ANTIC.
    \item Situer le règlement COBAC R-2016/04, la Convention de Malabo et
          la loi sur les données personnelles.
    \item Faire de la loi une source de critères d'audit.
    \item Appliquer les principes déontologiques de l'auditeur et
          reconnaître les menaces sur son indépendance.
    \item Rédiger un constat sans qualifier juridiquement les faits.
\end{itemize}
\end{boxobjectifs}

\bigskip

% ============================================================================
\section*{Avant de commencer : trois lettres}
% ============================================================================

Trois lettres, la même semaine de 2028\footnote{Les organisations, les
personnes, les dates et les montants des cas sont fictifs~; les textes
cités sont exacts, avec le degré de vérification indiqué.}.

À Yaoundé, M.~ESSOMBA reçoit de l'Agence nationale des technologies de
l'information et de la communication (\sigle{ANTIC}) l'annonce d'un audit
de sécurité de l'ANPP. Il se demande si son agence y est vraiment soumise,
et ce qu'il doit remettre.

À Douala, M.~ETEME écrit à AUDIX CONSEIL~: la synthèse destinée à
PHARMACAM « pourrait être formulée plus positivement », d'autant que le
renouvellement du contrat de LiZbiZ, et donc la suite des missions
d'AUDIX, en dépend.

À la BANQUE SANAGA enfin, la DSI propose au chef de mission de CERTIQUAL
AFRIQUE, la veille de l'étape~2, un « petit-déjeuner de travail » dans un
grand hôtel.

\medskip

Ce chapitre présente les textes qui encadrent la sécurité et son audit au
Cameroun et en Afrique, puis les règles de déontologie qui protègent
l'auditeur et la valeur de son travail.

\begin{boxalert}
\textbf{Règle du chapitre.} Un texte de loi n'est cité entre guillemets
que s'il a été lu dans sa version officielle. Les autres dispositions sont
présentées sans citation, avec la mention « à vérifier ». Un auditeur
applique la même règle dans ses rapports~: une citation inexacte ruine la
crédibilité de tout le reste.
\end{boxalert}

\separateur

% ============================================================================
\section{Panorama des Textes}
\label{sec:c10-panorama}
% ============================================================================

\begin{tableaucours}{L{3.2cm}L{3.0cm}L{3.8cm}L{2.2cm}}{%
    \tableauHeader{Texte} &
    \tableauHeader{Assujettis} &
    \tableauHeader{Ce qu'il impose} &
    \tableauHeader{Vérification}
}
Loi n°~2010/012 (cybersécurité et cybercriminalité) & Opérateurs,
exploitants de SI, autorités de certification, organismes publics & Mesures
de sécurité, gestion des risques, conservation des journaux, audit de
sécurité & Texte officiel lu \\
Règlement COBAC R-2016/04 & Établissements de crédit de la CEMAC &
Système de contrôle interne, dont le risque opérationnel & Art.~1 et~2
relevés, à collationner \\
Convention de Malabo (Union africaine, 2014) & États parties & Cadres
nationaux de cybersécurité et de protection des données & Statut
camerounais à vérifier \\
Loi n°~2024/017 (données personnelles) & Responsables de traitement &
Protection des données, y compris pendant l'audit & À vérifier au Journal
officiel \\
Droit OHADA des sociétés & Sociétés commerciales & Commissaire aux
comptes, comité d'audit & Présenté, à vérifier \\
\end{tableaucours}
\vspace{-0.8em}
\begin{center}{\small\itshape Tableau~10.1. Les textes et leur degré de
vérification dans ce cours}\end{center}

% ============================================================================
\section{La Loi sur la Cybersécurité}
\label{sec:c10-2010}
% ============================================================================

\subsection{L'audit de sécurité et l'ANTIC}

Le Chapitre~\ref{chap:audit} a cité la définition légale de l'audit de
sécurité (article~4, point~9) et le régime d'audit obligatoire et
périodique (articles~13 et~32)~\cite{loi2010012}. La loi confie à l'ANTIC,
désignée comme «~l'Agence~», une mission de contrôle de ces audits.

\begin{boxjuris}
\textbf{\loicam{n°~2010/012}, article~7, alinéa~2} (extrait)~: l'Agence a
notamment pour mission \textit{« de s'assurer de la régularité, de
l'effectivité des audits de sécurité des systèmes d'information suivant
les normes en la matière, des organismes publics et des autorités de
certification »}.

\medskip
\textbf{Portée.} Deux points. La loi renvoie aux \emph{normes}~: un audit
de sécurité légal se conduit selon les normes et référentiels des
Chapitres~\ref{chap:audit} et~\ref{chap:iso27000}. Et elle vise
expressément les \emph{organismes publics}, ce qui concerne l'ANPP.
\end{boxjuris}

\subsection{L'ANPP est-elle soumise~?}

\begin{boxscenario}{Opération ANPP. Lire le texte avant de conclure}
\textbf{Constatations.}
\begin{itemize}
    \item L'article~7, alinéa~2, charge l'Agence de s'assurer de la
          régularité des audits de sécurité des organismes publics.
    \item L'article~13, alinéa~1, soumet à l'audit obligatoire «~les
          réseaux de communications électroniques et les systèmes
          d'information, des opérateurs, des autorités de certification et
          des fournisseurs de services de communications électroniques~».
    \item Les conditions et modalités sont renvoyées à des textes
          d'application.
\end{itemize}

\textbf{Analyse} (confiance faible sur la portée de l'article~13, élevée
sur l'article~7). La ponctuation de l'article~13 admet deux lectures~:
soit tous les systèmes d'information sont visés, soit seulement ceux des
opérateurs, autorités et fournisseurs. L'article~7 tranche en partie pour
l'ANPP~: ses audits de sécurité relèvent du contrôle de l'Agence en tant
qu'organisme public. La réponse sûre à M.~ESSOMBA est donc~: l'audit
annoncé s'inscrit dans la mission légale de l'ANTIC~; ses modalités se
lisent dans les textes d'application en vigueur, que l'ANPP doit demander.
Un auditeur ne tranche pas seul une ambiguïté de texte~: il la signale.
\end{boxscenario}

\subsection{Des obligations qui deviennent des critères}

Plusieurs articles imposent aux exploitants des systèmes d'information des
obligations qui deviennent, pour l'auditeur, des critères opposables. Le
tableau les résume d'après le texte officiel~; seul l'article~29, alinéa~1,
a été cité textuellement (Chapitre~\ref{chap:iso27005}).

\begin{tableaucours}{L{1.6cm}L{6.2cm}L{4.4cm}}{%
    \tableauHeader{Art.} &
    \tableauHeader{Obligation (résumé du texte officiel)} &
    \tableauHeader{Critère d'audit}
}
26 (1) & Prendre les mesures techniques et administratives de sécurité~;
se doter de systèmes normalisés de gestion continue des risques & Un
processus de gestion des risques existe et fonctionne \\
26 (2) et (3) & Mettre en place des mécanismes de sécurité (disponibilité,
intégrité, authentification, non-répudiation, confidentialité, sécurité
physique), approuvés et visés par l'Agence & Les mécanismes existent~; le
visa de l'Agence est au dossier \\
29 (1) & Conserver les données de connexion et de trafic pendant dix ans &
La politique de conservation couvre dix ans \\
29 (2) & Installer des mécanismes de surveillance du contrôle d'accès aux
données & Les accès aux données sont journalisés et surveillés \\
30 (1) & Évaluer et réviser les systèmes de sécurité selon l'évolution des
technologies & Une revue périodique est tracée \\
\end{tableaucours}
\vspace{-0.8em}
\begin{center}{\small\itshape Tableau~10.2. Obligations de la loi
n°~2010/012 et critères d'audit}\end{center}

\begin{boxlecture}{L'ANPP face à la loi}
\textbf{Constatation.} À partir de la base commune du
Chapitre~\ref{chap:referentiels}, l'équipe de l'ANPP cote les cinq
obligations~:
\begin{itemize}
    \item 26~(1)~: \NON{}~: aucun registre des risques, aucune appréciation
          formalisée~;
    \item 26~(2) et~(3)~: \NEK{}~: des mécanismes existent, mais personne ne
          sait s'ils ont été soumis au visa de l'Agence~;
    \item 29~(1)~: \NEK{}~: les journaux sont «~aux archives~», leur durée
          de conservation n'est pas connue~;
    \item 29~(2)~: \NON{}~: les comptes des agents partis restent actifs, et
          aucune surveillance des accès n'est en place~;
    \item 30~(1)~: \NON{}~: aucune revue de sécurité depuis 2023.
\end{itemize}
0~\OUI{}, 3~\NON{}, 2~\NEK{}~: $\mu = 0$~; $U = 0{,}40$~;
$H = -(0{,}6\log_2 0{,}6 + 0{,}4\log_2 0{,}4) \approx 0{,}97$~bit.

\medskip
\textbf{Analyse} (confiance moyenne). Non-conformité opaque~: aucune
obligation n'est démontrée. Les deux \NEK{} se lèvent à peu de frais (une
demande à l'ANTIC sur le visa, une extraction des archives de journaux)~;
l'ANPP a intérêt à le faire \emph{avant} l'audit, pour que l'Agence
constate des écarts connus et traités plutôt qu'une ignorance.

\medskip
\hfill\lecture{0{,}00}{0{,}40}{0{,}97}
\end{boxlecture}

\subsection{Le secret, et l'obstruction}

\begin{boxjuris}
\textbf{\loicam{n°~2010/012}, article~14}~: \textit{« Le personnel de
l'Agence et les experts commis en vue d'accomplir des opérations d'audit
sont astreints au secret professionnel. »}

\medskip
\textbf{Article~61, alinéa~1}~: \textit{« Sont punis d'un emprisonnement de
trois (03) mois à trois (03) ans et d'une amende de 20.000 (vingt mille) à
100.000 (cent mille) F CFA, les personnels de l'Agence et les experts des
personnes morales chargés des audits qui révèlent sans autorisation, des
informations confidentielles dont ils ont eu connaissance à l'occasion
d'un audit de sécurité. »}

\medskip
\textbf{Article~61, alinéa~3} (paraphrasé)~: la loi punit aussi celui qui
fait obstacle aux audits de sécurité ou refuse de fournir les informations
ou documents qui s'y rapportent.

\medskip
\textbf{Portée.} La loi protège l'audit dans les deux sens~: elle punit
l'auditeur qui divulgue, et l'audité qui fait obstruction. Dans un audit de
sécurité légal, le refus du prestataire des GAB (Chapitre~\ref{chap:audit})
ne serait pas seulement une limitation de l'étendue~: ce serait un fait à
consigner, que l'Agence apprécierait.
\end{boxjuris}

% ============================================================================
\section{La Réglementation Bancaire}
\label{sec:c10-cobac}
% ============================================================================

Les établissements de crédit de la CEMAC sont soumis au règlement COBAC
R-2016/04 relatif au contrôle interne~\cite{cobac2016r04}.

\begin{boxjuris}
\textbf{Règlement COBAC R-2016/04, articles~1 et~2} (texte relevé sur une
base réglementaire en ligne, à collationner avec la version du
Secrétariat général de la COBAC)~:

\medskip
Article~1~: \textit{« Le présent règlement fixe les exigences minimales
relatives au système de contrôle interne que les établissements assujettis
doivent mettre en œuvre. »}

\medskip
Article~2, risque opérationnel~: \textit{« risque de pertes résultant de
l'inadéquation ou de la défaillance des processus internes, des personnes
et des systèmes, ou d'évènements externes. »}

\medskip
\textbf{Portée.} Une attaque par rançongiciel est un risque opérationnel
au sens du règlement~: défaillance des systèmes et des personnes,
provoquée par un événement externe. Le rapport demandé par la COBAC après
l'incident de septembre 2026 relève donc du contrôle interne de la banque,
et non d'une simple information technique.
\end{boxjuris}

Le règlement donne le cadre~; il ne dit pas comment évaluer une mesure de
sécurité. Pour le rapport au superviseur, la banque s'appuie sur ses
propres critères (ISO/IEC~27002, plan de traitement) et sur les résultats
de son audit interne. Les exigences précises du règlement sur les systèmes
d'information, au-delà des deux articles relevés, sont à lire dans le
texte collationné.

% ============================================================================
\section{Le Cadre Africain et les Données Personnelles}
\label{sec:c10-afrique}
% ============================================================================

\subsection{La Convention de Malabo}

La Convention de l'Union africaine sur la cybersécurité et la protection
des données à caractère personnel a été adoptée à Malabo en
2014~\cite{ua2014malabo}. Elle invite les États parties à se doter d'un
cadre national de cybersécurité, d'une protection des données personnelles
et d'une répression de la cybercriminalité. Elle est entrée en vigueur en
2023, après avoir réuni le nombre de ratifications requis. L'état de sa
signature et de sa ratification par le Cameroun est \emph{à vérifier}~;
le cours ne tire donc d'elle aucune obligation directe pour les
organisations camerounaises. Elle reste une référence pour lire la loi
nationale et pour les coopérations régionales.

\subsection{Les données personnelles pendant l'audit}

La loi n°~2024/017 relative à la protection des données à caractère
personnel est étudiée dans l'UE MSN-2S3-JUR01~; son texte reste à vérifier
au Journal officiel. Pour l'auditeur, une règle de prudence suffit~:
lorsqu'il extrait des données personnelles (fichiers du personnel,
bénéficiaires de paiements, clients), il ne prélève que ce que le test
exige, les protège pendant la mission, et les détruit au terme de la
conservation du dossier. L'audit de protection des données lui-même relève
de l'UE MSN-2S4-SEC02.

% ============================================================================
\section{La Déontologie de l'Auditeur}
\label{sec:c10-deontologie}
% ============================================================================

\subsection{Cinq principes}

Les normes mondiales d'audit interne de l'IIA regroupent l'éthique de
l'auditeur en cinq principes~\cite{iia2024gias}~; les codes de l'ISACA et
des organismes de certification en reprennent l'esprit.

\begin{tableaucours}{L{3.2cm}L{9.0cm}}{%
    \tableauHeader{Principe} &
    \tableauHeader{En audit de cybersécurité}
}
Intégrité & Ne rien affirmer que la preuve ne soutient pas~; ne rien
retirer sous la pression \\
Objectivité & Juger sans conflit d'intérêts ni préjugé~; ne pas auditer son
propre travail \\
Compétence & N'accepter que les missions que l'on sait conduire, ou se
faire assister (auditeur technique) \\
Conscience professionnelle & Le soin qu'exigent les enjeux~: tests
suffisants, papiers de travail complets \\
Confidentialité & Protéger les failles constatées et les données
consultées~; ne rien utiliser à des fins personnelles \\
\end{tableaucours}
\vspace{-0.8em}
\begin{center}{\small\itshape Tableau~10.3. Les principes déontologiques
(d'après les normes de l'IIA, 2024)}\end{center}

\subsection{Les menaces sur l'indépendance}

Le cadre conceptuel du code de déontologie de l'IESBA distingue cinq
menaces qui peuvent compromettre l'objectivité d'un
professionnel~\cite{iesba2023code}. Les cas du cours en offrent une
illustration chacun.

\begin{tableaucours}{L{2.6cm}L{4.6cm}L{5.0cm}}{%
    \tableauHeader{Menace} &
    \tableauHeader{Exemple du cours} &
    \tableauHeader{Sauvegarde}
}
Intérêt personnel & AUDIX dépend du contrat de LiZbiZ avec PHARMACAM pour
ses missions futures & Refuser toute modification non justifiée~;
déclarer la dépendance~; revue par un associé \\
Autorévision & Animer l'étude EBIOS que l'on auditera ensuite
(Chapitre~\ref{chap:ebios}) & Refuser l'une des deux missions \\
Représentation & Défendre l'audité devant son client ou son superviseur &
Ne s'exprimer que sur les constats~; renvoyer à l'audité \\
Familiarité & Auditer pendant des années la même équipe, devenue proche &
Rotation des chefs de mission \\
Intimidation & Une direction qui menace de retirer la mission si un
constat est maintenu & Remonter au commanditaire~; maintenir le constat \\
\end{tableaucours}
\vspace{-0.8em}
\begin{center}{\small\itshape Tableau~10.4. Les cinq menaces et leurs
sauvegardes (d'après le code de l'IESBA)}\end{center}

\begin{boxscenario}{Trois lettres, trois réponses}
\textbf{À M.~ETEME.} M.~BEYALA répond que la synthèse reprend les constats
tels qu'ils ont survécu au contradictoire~; il peut en améliorer la
clarté, mais non le sens. Il signale à son associé la dépendance de son
cabinet envers ce contrat (intérêt personnel, intimidation).

\medskip
\textbf{À la DSI de la BANQUE SANAGA.} Le chef de mission de CERTIQUAL
AFRIQUE décline l'invitation, et les règles de son organisme l'y obligent~:
un avantage offert par l'audité, même modeste, crée une apparence de
dépendance. Le refus figure au dossier.

\medskip
\textbf{À M.~ESSOMBA.} Les auditeurs de l'ANTIC sont astreints au secret
(article~14)~; l'ANPP, de son côté, doit fournir les documents demandés
(article~61, alinéa~3). La relation n'est ni hostile ni complaisante~:
elle est encadrée.
\end{boxscenario}

\subsection{Constater n'est pas qualifier}

\begin{boxalert}
\textbf{Violation typique.} Écrire dans un rapport « la DSI a dissimulé
l'absence de journalisation » ou « l'ANPP viole la loi ».

\smallskip
\textbf{Règle~:} l'auditeur constate un écart à un critère, preuve à
l'appui~: « pendant la période auditée, 7 agences sur 38 n'envoyaient pas
de journaux au collecteur~; le rapport de la RSSI de février 2028 les
déclarait raccordées ». L'intention (dissimuler) et la qualification
juridique (violation, infraction) ne se constatent pas~: la première
relève d'une enquête, la seconde du juge ou de l'autorité compétente.
\end{boxalert}

\subsection{Une obligation de moyens}

L'auditeur contractuel s'engage à conduire ses travaux selon les normes
qu'il déclare dans sa lettre de mission. Il n'est pas garant de l'absence
de faille~: son engagement est, en règle générale, une obligation de
moyens. Sa responsabilité peut être recherchée s'il n'a pas appliqué ces
normes. D'où l'importance des papiers de travail (Chapitre~\ref{chap:audit}),
qui montrent ce qu'il a fait, et des \NEK{} documentés, qui montrent ce
qu'il n'a pas pu établir~: ils le protègent autant qu'ils informent le
commanditaire.

% ============================================================================
\section{Ce que la Loi et la Déontologie demandent à chaque Métier}
\label{sec:c10-metiers}
% ============================================================================

\begin{boxmetier}{Auditeur en cybersécurité}
\textbf{Protéger la mission}
\begin{enumerate}
    \item Établir si la mission relève d'un régime légal, et en tirer les
          obligations de forme, de secret et de destinataire.
    \item Ne citer que des textes lus dans leur version officielle.
    \item Identifier, avant la mission, les menaces sur son indépendance,
          et les sauvegardes.
    \item Constater sans qualifier~; maintenir les constats sous la
          pression, et consigner la pression.
\end{enumerate}
\end{boxmetier}

\begin{boxmetier}{RSSI}
\textbf{Préparer l'audit légal}

Établir si l'organisation est assujettie, rassembler les pièces des
articles~26, 29 et~30 (appréciation des risques, mécanismes et visa,
conservation, revues), lever à l'avance les \NEK{} qui peuvent l'être, et
désigner l'interlocuteur des auditeurs~: l'obstruction est une infraction.
\end{boxmetier}

\begin{boxmetier}{Responsable GRC}
\textbf{Un registre des obligations}

Chaque article qui crée une obligation devient une ligne d'un registre des
obligations, reliée aux contrôles de la base commune
(Chapitre~\ref{chap:referentiels}) et à sa preuve. Le
Chapitre~\ref{chap:grc-pratique} montre comment le tenir.
\end{boxmetier}

\begin{boxmetier}{Consultant en mise en œuvre de SMSI}
\textbf{Conseiller sans auditer}

Le consultant qui a construit le SMSI ne participe pas à son audit, ni à
celui de l'organisme certificateur pour le compte de ce dernier. S'il est
sollicité, il le refuse ou le déclare.
\end{boxmetier}

\separateur

\begin{boxresume}
\textbf{Ce qu'il faut retenir de ce chapitre~:}
\begin{itemize}
    \item \textbf{Ne citer que le vérifié}~: chaque texte a son degré de
          vérification, et le rapport le dit.
    \item \textbf{La loi n°~2010/012} institue l'audit de sécurité, confie
          à l'ANTIC le contrôle des audits des organismes publics et des
          autorités de certification, et fournit des critères (articles~26,
          29, 30).
    \item \textbf{Elle protège l'audit dans les deux sens}~: secret de
          l'auditeur (articles~14 et~61), sanction de l'obstruction.
    \item \textbf{Le règlement COBAC R-2016/04} fait du risque informatique
          un risque opérationnel du contrôle interne.
    \item \textbf{La Convention de Malabo} encadre l'action des États~; son
          statut au Cameroun est à vérifier.
    \item \textbf{Cinq principes, cinq menaces}~: intégrité, objectivité,
          compétence, conscience professionnelle, confidentialité~; intérêt
          personnel, autorévision, représentation, familiarité,
          intimidation.
    \item \textbf{Constater n'est pas qualifier.}
\end{itemize}
\end{boxresume}

\begin{boxinfo}
\textbf{Points de contrôle.} Vous avez compris ce chapitre si vous savez~:
\begin{itemize}
    \item dire quel texte s'applique à une situation, et avec quel degré de
          certitude~;
    \item transformer un article de loi en critère d'audit~;
    \item nommer la menace sur l'indépendance dans une situation donnée, et
          la sauvegarde~;
    \item récrire une phrase de rapport qui qualifie au lieu de constater.
\end{itemize}
\end{boxinfo}

% ============================================================================
\section*{Exercices}
% ============================================================================

\begin{boxexo}[title={\ding{45}\quad Exercice~10.1. Quel texte~? \niveaubase}]
Pour chaque situation, indiquez le texte principal applicable et son degré
de vérification dans le cours.
\begin{enumerate}[(a)]
    \item La BANQUE SANAGA doit rendre compte au superviseur de l'incident
          de septembre 2026.
    \item Les journaux du pare-feu sont effacés au bout de trente jours.
    \item Un auditeur de l'ANTIC parle à un journaliste d'une faille de
          l'ANPP.
    \item Le prestataire des GAB refuse de remettre ses configurations lors
          d'un audit de sécurité légal.
    \item L'auditeur extrait le fichier du personnel de l'ANPP pour un
          rapprochement.
    \item Un étudiant affirme que la Convention de Malabo oblige LiZbiZ à
          notifier ses incidents.
\end{enumerate}
\end{boxexo}

\begin{boxexo}[title={\ding{45}\quad Exercice~10.2. Menaces sur l'indépendance \niveaumoyen}]
Nommez la menace et proposez une sauvegarde.
\begin{enumerate}[(a)]
    \item Mme~EKAMBI audite le plan de traitement qu'elle a aidé à rédiger
          en 2026.
    \item Le même chef de mission de CERTIQUAL AFRIQUE audite la banque
          depuis six ans.
    \item Un cabinet d'audit propose à LiZbiZ, pendant sa mission, de
          vendre l'outil qui corrigerait les constats.
    \item Le directeur général d'une agence publique fait savoir au
          magistrat qu'il «~en parlera au ministre~» si un constat est
          maintenu.
    \item Un auditeur témoigne devant la COBAC en faveur de la banque qu'il
          a auditée.
\end{enumerate}
\end{boxexo}

\begin{boxexo}[title={\ding{45}\quad Exercice~10.3. Constater, pas qualifier \niveaumoyen}]
Récrivez chacune des phrases suivantes sous la forme d'un constat,
preuve à l'appui, sans qualification juridique ni jugement d'intention.
\begin{enumerate}[(a)]
    \item « L'ANPP viole l'article~29 de la loi. »
    \item « Le responsable informatique de LiZbiZ a laissé des portes
          ouvertes aux fraudeurs. »
    \item « La banque a menti au certificateur dans sa déclaration
          d'applicabilité. »
    \item « Les quatre agents payés après leur départ ont détourné des
          fonds publics. »
\end{enumerate}
\end{boxexo}

\begin{boxexo}[title={\ding{45}\quad Exercice~10.4. Préparer l'audit de l'ANTIC \niveauexpert}]
M.~ESSOMBA vous confie la préparation de l'audit annoncé.
\begin{enumerate}[(a)]
    \item Rédigez la note qui répond à sa question «~sommes-nous soumis~?~»,
          en distinguant ce que le texte établit et ce qui reste à vérifier.
    \item Proposez, pour chacun des deux \NEK{} du Tableau~10.2 appliqué à
          l'ANPP, l'action qui le lève avant l'audit, et recalculez $\mu$,
          $U$ et $H$ si les deux deviennent \NON{}. La situation est-elle
          meilleure~? Justifiez.
    \item Listez les pièces à remettre aux auditeurs, et les précautions à
          prendre pour les données personnelles qu'elles contiennent.
    \item Que doit faire l'ANPP si un prestataire refuse de communiquer un
          document demandé par l'Agence~?
\end{enumerate}
\end{boxexo}

\bigskip

\begin{boxinfo}
\textbf{Pour aller plus loin.}
\begin{itemize}
    \item La loi n°~2010/012~\cite{loi2010012}, dans son texte officiel, et
          ses textes d'application.
    \item Le règlement COBAC R-2016/04~\cite{cobac2016r04}, dans la version
          du Secrétariat général de la COBAC.
    \item La Convention de Malabo~\cite{ua2014malabo}.
    \item Les normes mondiales d'audit interne~\cite{iia2024gias} et le code
          de l'IESBA~\cite{iesba2023code}.
\end{itemize}
\end{boxinfo}

\begin{boxinfo}
\textbf{Transition.} La Partie~III a posé les normes, les référentiels et
le droit. La Partie~IV passe à la pratique~: mettre en œuvre et auditer un
SMSI sur toute la BANQUE SANAGA, auditer la technique, faire vivre la GRC.
Elle s'ouvre avec le Chapitre~\ref{chap:smsi}.
\end{boxinfo}
